\section{模型假设}

题定的计算图、操作周期、缓存容量、总带宽及场景等待规则均作为固定输入；最终时长由官方确定性模拟器给出，不把这些题定条件另称为假设。为把不可直接枚举的切图与调度组合收敛到可计算的候选集合，本文只在构造方案时采用以下工作假设。

其一，把非 COPY 操作的给定计算周期加一个固定单位作为连通块的粗略负载权重。该权重只用于决定连通块向核心的初始分配，不预言实际运行时间；加一避免零周期操作在排序中完全失重。由于不同流水线可重叠、DDR 搬运会争用，这一代理可能低估通信密集块的代价，因而最终选择必须重新使用官方全局评估。

其二，把收缩 COPY 后互无依赖的弱连通块视为可独立分配的基本单元。这里的“独立”仅指该收缩图中没有非 COPY 操作依赖，并不假定它们在共享 DDR 上互不干扰。对于存在较大连通块的图，允许再沿拓扑深度切分；由此产生的跨段边界须由官方规则插入搬运并验证无环。

其三，候选搜索采用确定性排序与固定的规模阈值，同分时按操作编号或核心编号打破并列。这保证相同输入重复运行得到相同方案。候选集合不是全局最优解的证明；保留多个不同粒度的方案并按官方 Makespan 比较，同时用额外搬运量处理同等时长的取舍。共享 L2 研究还需把同一方案在有、无 L2 两种条件下成对评估，避免把方案变化误当作 Cache 效益。
